\section{Разделимые и префиксные схемы кодирования. Неравенство Макмиллана. Кодирование с минимальной избыточностью. Цена кодирования. Примеры.}

\subsection*{Схемы кодирования}

Пусть $M = \{m_1, \dots, m_n\}$ — алфавит сообщений, а $A = \{a_1, \dots, a_d\}$ — кодирующий алфавит. Схема алфавитного кодирования сопоставляет каждому символу $m_i$ слово $w_i$ в алфавите $A$.

\paragraph{Разделимые коды}
Схема кодирования называется \textbf{разделимой} (однозначно декодируемой), если любое слово в алфавите $A$, полученное в результате кодирования некоторой последовательности сообщений, может быть расшифровано единственным способом.
\textit{Пример:} $\{0, 01\}$ — неразделим (непонятно, $01$ — это одно сообщение или два: $0$ и $1$).

\paragraph{Префиксные коды (Условие Фано)}
Схема кодирования называется \textbf{префиксной}, если ни одно кодовое слово не является префиксом (началом) другого кодового слова.
\begin{itemize}
    \item \textbf{Свойство:} Любой префиксный код является разделимым.
    \item \textbf{Преимущество:} Декодирование может осуществляться мгновенно (нам не нужно ждать конца сообщения, чтобы понять, где закончилось текущее кодовое слово).
\end{itemize}
\textit{Пример префиксного кода:} $\{0, 10, 110, 111\}$.

\subsection*{Неравенство Макмиллана}

Эта теорема устанавливает необходимое и достаточное условие существования разделимого кода с заданными длинами слов $l_1, l_2, \dots, l_n$.

\paragraph{Теорема}
Для того чтобы существовал разделимый код в алфавите мощности $d$ с длинами кодовых слов $l_1, l_2, \dots, l_n$, необходимо и достаточно выполнение неравенства:
$$\sum_{i=1}^n d^{-l_i} \le 1.$$
\textit{Важное следствие:} Для любого разделимого кода существует префиксный код с такими же длинами слов. Поэтому на практике ограничиваются поиском префиксных кодов.

\subsection*{Цена и избыточность кодирования}

\paragraph{Цена кодирования (средняя длина)}
Если сообщения $m_i$ появляются с вероятностями $p_i$, то средней длиной (ценой) кода называется математическое ожидание длины кодового слова:
$$L = \sum_{i=1}^n p_i \cdot l_i.$$

\paragraph{Кодирование с минимальной избыточностью}
Задача состоит в поиске такой префиксной схемы, при которой цена $L$ минимальна для данных вероятностей $p_i$.
\begin{itemize}
    \item Теоретический предел минимальной цены задается энтропией источника $H$.
    \item \textbf{Избыточность} — это разность между средней длиной кода $L$ и энтропией $H$.
\end{itemize}

\subsection*{Алгоритм Хаффмана (Пример)}
Это жадный алгоритм построения оптимального префиксного кода.
1. Выбираются два сообщения с наименьшими вероятностями.
2. Они объединяются в новый узел с суммарной вероятностью.
3. Процесс повторяется, пока не останется один узел (корень дерева).
4. Путь от корня до листа определяет код (0 — поворот влево, 1 — вправо).

\paragraph{Пример}
Сообщения $A(0.5), B(0.25), C(0.25)$.
- Объединяем $B$ и $C$ (сумма $0.5$).
- Объединяем результат с $A$ (сумма $1.0$).
- Коды: $A=0, B=10, C=11$.
- Цена $L = 0.5\cdot 1 + 0.25\cdot 2 + 0.25\cdot 2 = 1.5$ бита на символ.